Ce livre rassemble les mathématiques enseignées en France de la sixième à la terminale, et
donne de chaque théorème du programme une démonstration vérifiée par une machine.

Chaque énoncé a été écrit dans \emph{Lean 4}, un assistant de preuve, puis démontré ; le
noyau de Lean a contrôlé chaque pas. Ce que vous lisez ici est la transcription française de
ces démonstrations : elle en suit les étapes une à une, sans les rendre plus élégantes.
Sous chacune, un lien mène à la ligne exacte du fichier vérifié.

Ce parti pris a une conséquence qu'il vaut mieux annoncer. Une démonstration formelle ne
saute rien. Là où un manuel écrit \og{} on montre de même \fg{}, il faut écrire le second
cas ; là où il écrit \og{} il est clair que \fg{}, il faut produire l'argument. Certaines
démonstrations paraîtront donc plus longues que dans un cours ordinaire. Elles ne sont pas
plus difficiles : elles sont complètes.

Une seconde conséquence est plus intéressante. Un programme scolaire \emph{admet} beaucoup :
le théorème des valeurs intermédiaires, l'existence des primitives, l'aire du disque. Un
assistant de preuve ne connaît pas cette convention — il faut démontrer, ou emprunter à une
bibliothèque qui a démontré. La frontière entre ce qui est admis et ce qui est établi se
déplace donc, et pas toujours dans le sens attendu : l'espérance d'une loi binomiale, admise
au lycée, est démontrée ici ; le théorème du toit, admis lui aussi, ne l'est pas. Chaque
chapitre dit où il s'arrête et pourquoi. Ces limites font partie du sujet.

Le lecteur n'a besoin de rien d'autre que du programme de sa classe. Aucune connaissance de
Lean n'est supposée : le code n'apparaît pas dans ces pages, seulement les mathématiques
qu'il contient.
